
%% Use the "normalphoto" option if you want a normal photo instead of cropped to a circle
% \documentclass[10pt,a4paper,normalphoto]{altacv}

\documentclass[10pt,a4paper,ragged2e,withhyper]{altacv}
%% AltaCV uses the fontawesome5 and simpleicons packages.
%% See http://texdoc.net/pkg/fontawesome5 and http://texdoc.net/pkg/simpleicons for full list of symbols.

% Change the page layout if you need to
\geometry{left=1.25cm,right=1.25cm,top=1.5cm,bottom=1.5cm,columnsep=1.2cm}

% The paracol package lets you typeset columns of text in parallel
\usepackage{paracol}

\usepackage{ragged2e}

% xelatex -shell-escape -output-driver="xdvipdfmx -z 0" sample.tex
\ifxetexorluatex
  % If using xelatex or lualatex:
  \setmainfont{Roboto Slab}
  \setsansfont{Lato}
  \renewcommand{\familydefault}{\sfdefault}
\else
  % If using pdflatex:
  \usepackage[rm]{roboto}
  \usepackage[defaultsans]{lato}
  % \usepackage{sourcesanspro}
  \renewcommand{\familydefault}{\sfdefault}
\fi

% Change the colours if you want to
\definecolor{SlateGrey}{HTML}{2E2E2E}
\definecolor{LightGrey}{HTML}{666666}
\definecolor{DarkPastelRed}{HTML}{450808}
\definecolor{PastelRed}{HTML}{8F0D0D}
\definecolor{GoldenEarth}{HTML}{E7D192}
\colorlet{name}{black}
\colorlet{tagline}{PastelRed}
\colorlet{heading}{DarkPastelRed}
\colorlet{headingrule}{GoldenEarth}
\colorlet{subheading}{PastelRed}
\colorlet{accent}{PastelRed}
\colorlet{emphasis}{SlateGrey}
\colorlet{body}{LightGrey}

% Change some fonts, if necessary
\renewcommand{\namefont}{\Huge\rmfamily\bfseries}
\renewcommand{\personalinfofont}{\footnotesize}
\renewcommand{\cvsectionfont}{\LARGE\rmfamily\bfseries}
\renewcommand{\cvsubsectionfont}{\large\bfseries}


% Change the bullets for itemize and rating marker
% for \cvskill if you want to
\renewcommand{\cvItemMarker}{{\small\textbullet}}
\renewcommand{\cvRatingMarker}{\faCircle}
% ...and the markers for the date/location for \cvevent
% \renewcommand{\cvDateMarker}{\faCalendar*[regular]}
% \renewcommand{\cvLocationMarker}{\faMapMarker*}


% If your CV/résumé is in a language other than English,
% then you probably want to change these so that when you
% copy-paste from the PDF or run pdftotext, the location
% and date marker icons for \cvevent will paste as correct
% translations. For example Spanish:
% \renewcommand{\locationname}{Ubicación}
% \renewcommand{\datename}{Fecha}


%% Use (and optionally edit if necessary) this .tex if you
%% want to use an author-year reference style like APA(6)
%% for your publication list
% \input{pubs-authoryear.tex}

%% Use (and optionally edit if necessary) this .tex if you
%% want an originally numerical reference style like IEEE
%% for your publication list
\input{pubs-num.tex}

%% sample.bib contains your publications
\addbibresource{sample.bib}

\begin{document}
\name{Lucas Petit}
\tagline{Available from mid-January 2025}
%% You can add multiple photos on the left or right
\photoR{3cm}{CV}
%\photoL{2.8cm}{hovercode}

\personalinfo{%
  % Not all of these are required!
  \email{lucas.noirotorion@gmail.com}
  \phone{+33685263667}
  %\mailaddress{Åddrésş, Street, 00000 Cóuntry}
  \location{Montpellier, France}
  \homepage{www.lucas-petit.fr}
  % \twitter{@twitterhandle}
  %\xtwitter{@x-handle}
  \linkedin{lucas-petit-dev}
  \github{Crowreed}
  %\orcid{0000-0000-0000-0000}
  %% You can add your own arbitrary detail with
  %% \printinfo{symbol}{detail}[optional hyperlink prefix]
  % \printinfo{\faPaw}{Hey ho!}[https://example.com/]

  %% Or you can declare your own field with
  %% \NewInfoFiled{fieldname}{symbol}[optional hyperlink prefix] and use it:
  % \NewInfoField{gitlab}{\faGitlab}[https://gitlab.com/]
  % \gitlab{your_id}
  %%
  %% For services and platforms like Mastodon where there isn't a
  %% straightforward relation between the user ID/nickname and the hyperlink,
  %% you can use \printinfo directly e.g.
  % \printinfo{\faMastodon}{@username@instace}[https://instance.url/@username]
  %% But if you absolutely want to create new dedicated info fields for
  %% such platforms, then use \NewInfoField* with a star:
  % \NewInfoField*{mastodon}{\faMastodon}
  %% then you can use \mastodon, with TWO arguments where the 2nd argument is
  %% the full hyperlink.
  % \mastodon{@username@instance}{https://instance.url/@username}
}

\makecvheader
%% Depending on your tastes, you may want to make fonts of itemize environments slightly smaller
% \AtBeginEnvironment{itemize}{\small}

%% Set the left/right column width ratio to 6:4.
\columnratio{0.6}

% Start a 2-column paracol. Both the left and right columns will automatically
% break across pages if things get too long.
\begin{paracol}{2}

%%%%%%%%%%%%%%%%%%%%%%%%%%%%%%%%%%%%%%%%%%%%%%%%%%%%%%%%%%%%%%%%

\cvsection{About me}
\begin{quote}
\begin{justify}
Passionate about mathematics and computer science, I find my balance in abstraction, where ideas take shape and become concrete solutions. I enjoy transforming theoretical complexity into tangible applications.
\end{justify}
\end{quote}

%%%%%%%%%%%%%%%%%%%%%%%%%%%%%%%%%%%%%%%%%%%%%%%%%%%%%%%%%%%%%%%%

\cvsection{Experience}

\cvevent{Private Tutor -- Mathematics}{Freelance}{}{Dijon}
\begin{itemize}
    \item Developed lesson plans tailored to individual needs.
    \item Simplified complex concepts to help students better understand abstract and technical subjects.
\end{itemize}

\divider

\cvevent{Production Operator}{Freudenberg}{}{Langres}
\begin{itemize}
\item Collaborated with a team to meet daily production targets within tight deadlines.
\end{itemize}

%%%%%%%%%%%%%%%%%%%%%%%%%%%%%%%%%%%%%%%%%%%%%%%%%%%%%%%%%%%%%%%%%

\cvsection{Projects}

\cvevent{Real Estate Price Prediction Project Using Linear Regression}{}{}{}

\begin{itemize}
    \item Implemented a linear regression model using Scikit-learn to predict house prices based on features from the Boston Housing dataset.
    \item Utilized Pandas for data cleaning and preparation, and Matplotlib to visualize relationships between variables.
    \item Optimized the model with cross-validation methods, achieving an $R^2$ score of $0.85$.
    
\end{itemize}



\cvevent{Design and implementation of ray tracing algorithms in Python.}{}{}{}

\divider

\cvevent{Development of a multiplayer game in C++ with the implementation of a communication protocol.}{}{}{}

\divider

\cvevent{Design and implementation of a heuristic for the WSRP problem in Java - hosted on GitLab.}{}{}{}
% use ONLY \newpage if you want to force a page break for
% ONLY the current column
\newpage
%% Switch to the right column. This will now automatically move to the second
%% page if the content is too long.
\switchcolumn

%%%%%%%%%%%%%%%%%%%%%%%%%%%%%%%%%%%%%%%%%%%%%%%%%%%%%%%%%%%%%%%%%

\cvsection{Skills}

\cvtag{Proactive}
\cvtag{Self-taught}
\cvtag{Ability to abstract}\\

\divider\smallskip

\cvtag{C, C++} \cvtag{HTML5, CSS3, JavaScript} \cvtag{Java} \cvtag{R}\\
\cvtag{Python, PyTorch, Pandas, Scikit-learn}

\divider\smallskip

\cvtag{Writing: LaTeX, Markdown}
\cvtag{Office Suite}\\
\cvtag{Linux, Windows}
\cvtag{GitHub, GitLab}

\divider\smallskip

\cvtag{Operations Research, MILP}\cvtag{Formal Computation}\cvtag{Software Security}

\divider\smallskip

\cvtag{Applied Mathematics}
\cvtag{Statistics and Probability}\\
\cvtag{Numerical Analysis}


%%%%%%%%%%%%%%%%%%%%%%%%%%%%%%%%%%%%%%%%%%%%%%%%%%%%%%%%%%%%%%%%%

\cvsection{Languages}

\large \textbf{French} - native

\large  \textbf{English} - C1 

%% Yeah I didn't spend too much time making all the
%% spacing consistent... sorry. Use \smallskip, \medskip,
%% \bigskip, \vspace etc to make adjustments.
\medskip

%%%%%%%%%%%%%%%%%%%%%%%%%%%%%%%%%%%%%%%%%%%%%%%%%%%%%%%%%%%%%%%%%

\cvsection{Education}

\cvevent{Master's Degree in Computer Science -- Specialization: Algorithms}{University of Montpellier}{September 2023 -- Present}{}
  \begin{itemize}
    \item \footnotesize \textbf{Master 1} - University of Montpellier
    \item \footnotesize \textbf{Master 2} - University of Montpellier (in progress)
  \end{itemize}

\divider

\cvevent{Master's Degree in Mathematics}{}{}{}
  \begin{itemize}
    \item \footnotesize \textbf{Master 1} - University of Burgundy
    \item \footnotesize \textbf{Master 2} - University of Montpellier (S1 completed)
  \end{itemize}
\divider

\cvevent{Bachelor's Degree in Mathematics}{University of Burgundy}{2018 -- 2021}{}

\divider

\cvevent{Scientific Baccalaureate -- Engineering Sciences}{Denis Diderot High School}{2015 -- 2018}{Langres}
\footnotesize Graduated with high honors, European distinction


\end{paracol}


\end{document}
